% Work-MAT | 2. rocnik | pripravovana hodina 33
% Návaznost: uživatelské připravy.zip (hodiny 01–20), autorská učebnice na GitHubu.
% Realisticky.cz: didaktická inspirace v posloupnosti a typech úloh; zadání jsou vlastní.
\documentclass[11pt,a4paper]{article}
\input{../../sablony/pripravy/priprava.sty}
\nastavPripravu{2. ročník}{Logaritmy}{33}{Základní hodnoty logaritmů}
\begin{document}
\blok{Motivace a návaznost}
Zopakuj krátkou tabulku mocnin pro základy 2, 3 a 10 a nech žáky číst logaritmy jako otázky: na kolikátou?
\poznamka{Orientační nabídka pro 45--55 minut: krátké zopakování a motivace, nový výklad, tři společné vzorové úlohy, samostatná práce. Vyber podle tempa třídy; zbývající zadání slouží jako rezerva.}
\blok{Výklad}
Platí $\log_a1=0$, $\log_a a=1$, $\log_a a^r=r$ pro přípustný základ. Dekadický logaritmus značíme $\log$ nebo $\log_{10}$, přirozený $\ln$.
\blok{Společné příklady}
\begin{prikladReseniBox}[title={Příklad 1}]
Vyřeš, resp. urči: $\log_2 32$.
\tcblower
\[\begin{aligned}2^5=32 \\ \log_2 32=5\end{aligned}\]
\end{prikladReseniBox}
\begin{prikladReseniBox}[title={Příklad 2}]
Vyřeš, resp. urči: $\log_{\frac{1}{2}}8$.
\tcblower
\[\begin{aligned}\left(\frac{1}{2}\right)^{-3}=8 \\ \log_{\frac{1}{2}}8=-3\end{aligned}\]
\end{prikladReseniBox}
\poznamka{U klesajícího základu mezi nulou a jednou nezaměňuj znaménko exponentu. Převody prověř i pro zlomky.}
\blok{Samostatné procvičování}
\begin{zadaniReseniBox}[title={Samostatná práce 1}]
\begin{enumerate}[label=\alph*)]
\item $\log_{2}\left(8\right)$.
\item $\log_{3}\left(9\right)$.
\item $\log_{5}\left(25\right)$.
\item $\log_{2}\left(32\right)$.
\item $\log_{4}\left(16\right)$.
\item $\log_{3}\left(27\right)$.
\end{enumerate}
\tcblower
\begin{enumerate}[label=\alph*)]
\item $3$.
\item $2$.
\item $2$.
\item $5$.
\item $2$.
\item $3$.
\end{enumerate}
\end{zadaniReseniBox}
\blok{Dobrovolné rozšíření}
\begin{zadaniReseniBox}[title={Rozšíření 1}]
\begin{enumerate}[label=\alph*)]
\item $Vytvoř vlastní úlohu tohoto typu a se spolužákem porovnej podmínky a výsledek.$.
\end{enumerate}
\tcblower
\begin{enumerate}[label=\alph*)]
\item Zkontroluj definice a argumenty..
\end{enumerate}
\end{zadaniReseniBox}
\blok{Úkol na doma}
U každého logaritmu vypiš podmínky argumentu a proveď zkoušku. U rovnic vyznač případně nepřípustné kořeny, u nerovnic odpověz intervalem.
\begin{zadaniReseniBox}[title={Domácí úkol}]
\begin{enumerate}[label=\alph*)]
\item $\log_{2}\left(\frac{1}{4}\right)$.
\item $\log_{5}\left(\frac{1}{5}\right)$.
\item $\log_{10}\left(1000\right)$.
\item $\log_{2}\left(1\right)$.
\item $\log_{3}\left(\frac{1}{9}\right)$.
\item $\log_{4}\left(64\right)$.
\item $\log_{5}\left(1\right)$.
\item $\log_{10}\left(\frac{1}{10}\right)$.
\end{enumerate}
\tcblower
\begin{enumerate}[label=\alph*)]
\item $-2$.
\item $-1$.
\item $3$.
\item $0$.
\item $-2$.
\item $3$.
\item $0$.
\item $-1$.
\end{enumerate}
\end{zadaniReseniBox}
\poznamka{Domácí úkol obsahuje 8 položek; předpokládaná domácí práce přibližně 15--30 minut. Po větších blocích jsou vědomě promíchané dříve probrané typy.}
\end{document}
